% El espaciado explícito de porcentajes se conserva también en modo matemático.
\AtBeginDocument{\spanishplainpercent}
% Declaración ajustada al trabajo real del Subdocumento 4.
\NewDocumentCommand{\FilaDeclaracionIA}{m m m m m m}{%
  #1 & #2 & #3 & #4 & #5 & #6 \\
  \hline
}

\providecommand{\DeclaracionesIA}{%
  \FilaDeclaracionIA{4 (intro.)}{OpenAI Codex}{Síntesis del alcance, dependencias y presentación.}{Alto}{Ninguno}{Fernando Guerrero: revisión manual actual de síntesis del alcance, dependencias y presentación y de su coherencia documental.}
  \FilaDeclaracionIA{4.1}{OpenAI Codex}{Apoyo a la redacción, contraste documental y coherencia de Arquitectura lógica.}{Alto}{Alto}{Fernando Guerrero: revisión manual actual de arquitectura lógica y de su coherencia documental. Revisión de los diagramas asociados.}
  \FilaDeclaracionIA{4.1.1}{OpenAI Codex}{Apoyo a la redacción, contraste documental y coherencia de Especificaciones Tecnologías de Software a utilizar.}{Alto}{Ninguno}{Fernando Guerrero: revisión manual actual de especificaciones Tecnologías de Software a utilizar y de su coherencia documental.}
  \FilaDeclaracionIA{4.2}{OpenAI Codex}{Apoyo a la redacción, contraste documental y coherencia de Arquitectura física.}{Alto}{Alto}{Fernando Guerrero: revisión manual actual de arquitectura física y de su coherencia documental. Revisión de los diagramas asociados.}
  \FilaDeclaracionIA{4.2.1}{OpenAI Codex}{Apoyo a la redacción, contraste documental y coherencia de Especificaciones Implementos a proveer (Hardware y Software).}{Alto}{Alto}{Fernando Guerrero: revisión manual actual de especificaciones Implementos a proveer (Hardware y Software) y de su coherencia documental. Revisión de los diagramas asociados.}
  \FilaDeclaracionIA{4.3}{OpenAI Codex}{Apoyo a la redacción, contraste documental y coherencia de Data center.}{Alto}{Alto}{Fernando Guerrero: revisión manual actual de data center y de su coherencia documental. Revisión de los diagramas asociados.}
  \FilaDeclaracionIA{4.3.1}{OpenAI Codex}{Apoyo a la redacción, contraste documental y coherencia de Especificaciones Data Center Primaria.}{Alto}{Alto}{Fernando Guerrero: revisión manual actual de especificaciones Data Center Primaria y de su coherencia documental. Revisión de los diagramas asociados.}
  \FilaDeclaracionIA{4.3.2}{OpenAI Codex}{Apoyo a la redacción, contraste documental y coherencia de Especificaciones Data Center Secundario.}{Alto}{Alto}{Fernando Guerrero: revisión manual actual de especificaciones Data Center Secundario y de su coherencia documental. Revisión de los diagramas asociados. José Lara: revisión de incoherencias entre la recuperación propuesta y los objetivos de continuidad de SD5 y SD8, incluyendo regiones, réplica y retorno.}
  \FilaDeclaracionIA{A4.1}{OpenAI Codex}{Registro de decisiones de arquitectura y alternativas.}{Alto}{Ninguno}{Fernando Guerrero: revisión manual actual de registro de decisiones de arquitectura y alternativas y de su coherencia documental.}
  \FilaDeclaracionIA{A4.2}{OpenAI Codex}{Contratos, cuotas, integraciones y dimensionamiento.}{Alto}{Ninguno}{Fernando Guerrero: revisión manual actual de contratos, cuotas, integraciones y dimensionamiento y de su coherencia documental.}
  \FilaDeclaracionIA{A4.3}{OpenAI Codex}{Memorias eléctricas, ambientales y de protección.}{Alto}{Medio}{Fernando Guerrero: revisión manual actual de memorias eléctricas, ambientales y de protección y de su coherencia documental. Revisión de los diagramas asociados.}
  \FilaDeclaracionIA{A4.4}{OpenAI Codex}{Implantación, plano exterior y gabinetes.}{Alto}{Alto}{Fernando Guerrero: revisión manual actual de implantación, plano exterior y gabinetes y de su coherencia documental.}
  \FilaDeclaracionIA{A4.5}{OpenAI Codex}{Configuración, gestión, energía y recepción de puestos.}{Alto}{Ninguno}{Fernando Guerrero: revisión manual actual de configuración, gestión, energía y recepción de puestos y de su coherencia documental.}
  \FilaDeclaracionIA{A4.6}{OpenAI Codex}{Capacidad, equivalencia de ambientes y cálculos.}{Alto}{Ninguno}{Fernando Guerrero: revisión manual actual de capacidad, equivalencia de ambientes y cálculos y de su coherencia documental.}
  \FilaDeclaracionIA{T-11}{OpenAI Codex}{Inventario, cantidades y concordancia con el diseño.}{Alto}{Ninguno}{Fernando Guerrero: revisión manual actual de inventario, cantidades y concordancia con el diseño y de su coherencia documental.}
}

\newenvironment{declaracionia}{%
  \begingroup
  \setlength{\tabcolsep}{3pt}
  \begin{longtable}{L{\dimexpr.105\textwidth-2\tabcolsep\relax}L{\dimexpr.14\textwidth-2\tabcolsep\relax}L{\dimexpr.21\textwidth-2\tabcolsep\relax}L{\dimexpr.115\textwidth-2\tabcolsep\relax}L{\dimexpr.14\textwidth-2\tabcolsep\relax}L{\dimexpr.29\textwidth-2\tabcolsep\relax}}
    \caption{Declaración de uso de IA por sección y archivo asociado}
    \label{tab:declaracion-ia} \\
    \hline
    \EncabezadoTabla{Sección} &
    \EncabezadoTabla{Herramienta} &
    \EncabezadoTabla{Finalidad del uso} &
    \EncabezadoTabla{Nivel en texto} &
    \EncabezadoTabla{Nivel en diagramas} &
    \EncabezadoTabla{Revisión humana: responsable y verificación} \\
    \hline
    \endfirsthead
    \hline
    \EncabezadoTabla{Sección} &
    \EncabezadoTabla{Herramienta} &
    \EncabezadoTabla{Finalidad del uso} &
    \EncabezadoTabla{Nivel en texto} &
    \EncabezadoTabla{Nivel en diagramas} &
    \EncabezadoTabla{Revisión humana: responsable y verificación} \\
    \hline
    \endhead
}{%
  \end{longtable}
  \endgroup
}

\newcommand{\CierreSubdocumento}{%
  \clearpage
  \printbibliography[heading=referenciassynaptix,title={Referencias}]
  \clearpage
  \chapter*{Declaración de uso de IA}
  \addcontentsline{toc}{chapter}{Declaración de uso de IA}
  \markboth{Declaración de uso de IA}{Declaración de uso de IA}
  Synaptix declara el uso de IA por sección, anexo y formulario asociado. La tabla identifica las herramientas, su finalidad, los niveles de intervención y quién realizó la revisión manual del contenido y de su coherencia documental. Se conservan los aportes previos de cada integrante. Synaptix asume la responsabilidad sobre el contenido; esta declaración se consolidará en el Formulario A-6.

  La Tabla~\ref{tab:declaracion-ia} presenta el detalle del uso declarado.
  \begin{declaracionia}
    \DeclaracionesIA
  \end{declaracionia}
  \FuenteTabla{Registro de uso de IA de este subdocumento.}

  El control previo a la emisión comprende modelos, certificaciones, distancias, disponibilidad regional, clasificación del surtidor, normas eléctricas y localización de las fuentes en las bases.
}
